% Options for packages loaded elsewhere
% Options for packages loaded elsewhere
\PassOptionsToPackage{unicode}{hyperref}
\PassOptionsToPackage{hyphens}{url}
\PassOptionsToPackage{dvipsnames,svgnames,x11names}{xcolor}
%
\documentclass[
  letterpaper,
  DIV=11,
  numbers=noendperiod]{scrartcl}
\usepackage{xcolor}
\usepackage[margin=1in]{geometry}
\usepackage{amsmath,amssymb}
\setcounter{secnumdepth}{-\maxdimen} % remove section numbering
\usepackage{iftex}
\ifPDFTeX
  \usepackage[T1]{fontenc}
  \usepackage[utf8]{inputenc}
  \usepackage{textcomp} % provide euro and other symbols
\else % if luatex or xetex
  \usepackage{unicode-math} % this also loads fontspec
  \defaultfontfeatures{Scale=MatchLowercase}
  \defaultfontfeatures[\rmfamily]{Ligatures=TeX,Scale=1}
\fi
\usepackage{lmodern}
\ifPDFTeX\else
  % xetex/luatex font selection
  \setmainfont[BoldFont=EB Garamond SemiBold,ItalicFont=EB Garamond
Italic,BoldItalicFont=EB Garamond SemiBold Italic]{EB Garamond Math}
  \setmathfont[]{EB Garamond Math}
\fi
% Use upquote if available, for straight quotes in verbatim environments
\IfFileExists{upquote.sty}{\usepackage{upquote}}{}
\IfFileExists{microtype.sty}{% use microtype if available
  \usepackage[]{microtype}
  \UseMicrotypeSet[protrusion]{basicmath} % disable protrusion for tt fonts
}{}
\makeatletter
\@ifundefined{KOMAClassName}{% if non-KOMA class
  \IfFileExists{parskip.sty}{%
    \usepackage{parskip}
  }{% else
    \setlength{\parindent}{0pt}
    \setlength{\parskip}{6pt plus 2pt minus 1pt}}
}{% if KOMA class
  \KOMAoptions{parskip=half}}
\makeatother
% Make \paragraph and \subparagraph free-standing
\makeatletter
\ifx\paragraph\undefined\else
  \let\oldparagraph\paragraph
  \renewcommand{\paragraph}{
    \@ifstar
      \xxxParagraphStar
      \xxxParagraphNoStar
  }
  \newcommand{\xxxParagraphStar}[1]{\oldparagraph*{#1}\mbox{}}
  \newcommand{\xxxParagraphNoStar}[1]{\oldparagraph{#1}\mbox{}}
\fi
\ifx\subparagraph\undefined\else
  \let\oldsubparagraph\subparagraph
  \renewcommand{\subparagraph}{
    \@ifstar
      \xxxSubParagraphStar
      \xxxSubParagraphNoStar
  }
  \newcommand{\xxxSubParagraphStar}[1]{\oldsubparagraph*{#1}\mbox{}}
  \newcommand{\xxxSubParagraphNoStar}[1]{\oldsubparagraph{#1}\mbox{}}
\fi
\makeatother


\usepackage{longtable,booktabs,array}
\newcounter{none} % for unnumbered tables
\usepackage{calc} % for calculating minipage widths
% Correct order of tables after \paragraph or \subparagraph
\usepackage{etoolbox}
\makeatletter
\patchcmd\longtable{\par}{\if@noskipsec\mbox{}\fi\par}{}{}
\makeatother
% Allow footnotes in longtable head/foot
\IfFileExists{footnotehyper.sty}{\usepackage{footnotehyper}}{\usepackage{footnote}}
\makesavenoteenv{longtable}
\usepackage{graphicx}
\makeatletter
\newsavebox\pandoc@box
\newcommand*\pandocbounded[1]{% scales image to fit in text height/width
  \sbox\pandoc@box{#1}%
  \Gscale@div\@tempa{\textheight}{\dimexpr\ht\pandoc@box+\dp\pandoc@box\relax}%
  \Gscale@div\@tempb{\linewidth}{\wd\pandoc@box}%
  \ifdim\@tempb\p@<\@tempa\p@\let\@tempa\@tempb\fi% select the smaller of both
  \ifdim\@tempa\p@<\p@\scalebox{\@tempa}{\usebox\pandoc@box}%
  \else\usebox{\pandoc@box}%
  \fi%
}
% Set default figure placement to htbp
\def\fps@figure{htbp}
\makeatother





\setlength{\emergencystretch}{3em} % prevent overfull lines

\providecommand{\tightlist}{%
  \setlength{\itemsep}{0pt}\setlength{\parskip}{0pt}}



 


\linespread{1.05}
\RedeclareSectionCommand[
  afterskip=0.3ex,
  beforeskip=0.8ex
  ]{subsection}
\RedeclareSectionCommand[
  afterskip=0.1ex,
  beforeskip=0.5ex
  ]{subsubsection}
\KOMAoption{captions}{tableheading}
\makeatletter
\@ifpackageloaded{caption}{}{\usepackage{caption}}
\AtBeginDocument{%
\ifdefined\contentsname
  \renewcommand*\contentsname{Table of contents}
\else
  \newcommand\contentsname{Table of contents}
\fi
\ifdefined\listfigurename
  \renewcommand*\listfigurename{List of Figures}
\else
  \newcommand\listfigurename{List of Figures}
\fi
\ifdefined\listtablename
  \renewcommand*\listtablename{List of Tables}
\else
  \newcommand\listtablename{List of Tables}
\fi
\ifdefined\figurename
  \renewcommand*\figurename{Figure}
\else
  \newcommand\figurename{Figure}
\fi
\ifdefined\tablename
  \renewcommand*\tablename{Table}
\else
  \newcommand\tablename{Table}
\fi
}
\@ifpackageloaded{float}{}{\usepackage{float}}
\floatstyle{ruled}
\@ifundefined{c@chapter}{\newfloat{codelisting}{h}{lop}}{\newfloat{codelisting}{h}{lop}[chapter]}
\floatname{codelisting}{Listing}
\newcommand*\listoflistings{\listof{codelisting}{List of Listings}}
\makeatother
\makeatletter
\makeatother
\makeatletter
\@ifpackageloaded{caption}{}{\usepackage{caption}}
\@ifpackageloaded{subcaption}{}{\usepackage{subcaption}}
\makeatother
\usepackage{bookmark}
\IfFileExists{xurl.sty}{\usepackage{xurl}}{} % add URL line breaks if available
\urlstyle{same}
% fallback for those not using the hyperref driver hyperxmp:
\makeatletter
\@ifundefined{xmpquote}{\newcommand{\xmpquote}[1]{#1}}{}
\makeatother
\hypersetup{
  pdftitle={PHIL 444: Game Theory and Social Choice},
  pdfauthor={Brian Weatherson},
  colorlinks=true,
  linkcolor={blue},
  filecolor={Maroon},
  citecolor={Blue},
  urlcolor={Blue},
  pdfcreator={LaTeX via pandoc}}


\title{PHIL 444: Game Theory and Social Choice}
\author{Brian Weatherson}
\date{Fall 2026}
\begin{document}
\maketitle


\textbf{Lead Instructor}: Brian Weatherson\\
\textbf{Email}: weath@umich.edu\\
\textbf{Web}: canvas.umich.edu\\
\textbf{Office Hours}: TBD

~

\section{Key Points}\label{key-points}

\begin{itemize}
\tightlist
\item
  This course examines two philosophical questions about choice in
  social settings.
\item
  The first half is on game theory, with a special focus on what we can
  rationally do when all we know is that everyone involved is rational,
  and what does this tell us about rationality?
\item
  The second half works through Amartya Sen's \emph{Collective Choice
  and Social Welfare} (2017 expanded edition): what does it mean for a
  group to choose, and how can individual judgments be combined into
  collective ones?
\item
  There will be more mathematics than in a typical philosophy course,
  but less than in an economics course covering this material. In
  particular, there will be no calculus.
\item
  Assessment consists of two essays (25\% each), eight quizzes with the
  best six counting for 40\% in total, and class participation (10\%).
\end{itemize}

\section{Course Description}\label{course-description}

This course examines two philosophical questions about choice in social
settings.

The first half asks what rational choice looks like when the unknowns
are not natural events (the weather, the bus arrival time) but the
actions of other rational agents who are reasoning about you. Standard
expected-utility theory does not pin down behavior in such settings. We
work through the central solution concepts of game theory (Nash
equilibrium, rationalizability, subgame perfection, perfect Bayesian
equilibrium) and ask what philosophical claims each one rests on. At the
end we discuss signaling games and refinements of equilibrium concepts,
the beer-quiche example and the Cho-Kreps intuitive criterion, and two
famous applications of these games: used car markets and job-market
signaling.

The second half asks what it means for a group to choose. Working from
Sen's \emph{Collective Choice and Social Welfare} (2017 expanded
edition), we study the formalism behind results like Arrow's famous
impossibility theorem and the philosophical responses Sen develops,
including his work on interpersonal comparison, justice, rights, and the
limits of welfarism.

The unifying theme is that individual rationality does not extend
cleanly to social settings. Game theory and social choice theory are two
related, but interestingly distinct, ways of seeing how social settings
create new complications. We will look at what philosophers, political
scientists, and economists, have made of these complications.

\section{Required Materials}\label{required-materials}

The required text is:

\begin{itemize}
\tightlist
\item
  Amartya Sen, \emph{Collective Choice and Social Welfare: An Expanded
  Edition} (Harvard University Press, 2017). Available electronically
  through the University of Michigan library. (Hereafter called CCSW.)
\end{itemize}

The main reading for the first half is my own lecture notes, \emph{Game
Theory and Social Choice}, free at
\url{https://bweatherson.github.io/F26-Phil444-site/notes/},
supplemented by three classic papers (Akerlof on lemons, Spence on
job-market signalling, Cho and Kreps on the intuitive criterion).
Giacomo Bonanno's \emph{Game Theory}, also free, is a good second voice
if you'd like something other than my lectures and my notes. Very
usefully, it comes with worked solutions to every exercise. It's
available at
\url{https://faculty.econ.ucdavis.edu/faculty/bonanno/PDF/GT_book.pdf}.
Section numbers in the schedule refer to that file, which is the third
edition. If you just Google it you might find earlier editions, which
have slightly different numbers. (ArXiV, for instance, just has the
second edition.)

There will be some supplementary readings in the second half, which will
be linked, but the priority is CCSW. All readings will be made available
through Canvas or the course website.

\section{Course Requirements}\label{course-requirements}

\begin{itemize}
\tightlist
\item
  There will be 2 essays, each worth 25\% of the grade, one for each
  half of the course.
\item
  There will be 8 quizzes over the semester. Your best 6 results count,
  and together they are worth 40\% of the grade.
\item
  10\% of the grade is for participation in class.
\end{itemize}

\section{Summary of Grading System}\label{summary-of-grading-system}

\begin{enumerate}
\def\labelenumi{\arabic{enumi}.}
\tightlist
\item
  Essays - 2 essays, 25\% each: 50\%
\item
  Quizzes - best 6 of 8: 40\%
\item
  Class participation: 10\%
\end{enumerate}

\section{Canvas}\label{canvas}

There is a Canvas site for this course, which can be accessed from
\url{https://canvas.umich.edu}. Course documents will be available from
this site. Please make sure that you can access this site. Consult the
site regularly for announcements, including changes to the course
schedule. And there are many tools on the site to communicate with each
other, and with me.

There is also a github site for the course which contains more
documents, and it is at
\url{https://bweatherson.github.io/F26-Phil444-site/}.

\pagebreak

\section{Class Schedule}\label{class-schedule}

Readings should be done before the scheduled class. Lectures 1-14 cover
game theory; Lectures 15-28 cover social choice theory, working through
Sen's \emph{Collective Choice and Social Welfare} (cited below as CCSW).
``Notes'' below refers to my lecture notes, \emph{Game Theory and Social
Choice}, at \url{https://bweatherson.github.io/F26-Phil444-site/notes/}.

CCSW alternates between chapters and starred chapters. The starred
chapters are more mathematical. Sen says in the preface that the
unstarred chapters can be read on their own, and we're \emph{mostly}
going to follow him. We will mostly be using the unstarred chapers as
reqiured reading, and the starred ones as recommended, and I'll go over
the key proofs. But we won't \emph{always} to this. In some cases the
formalism is the point of what we're studying, and in those cases we'll
be focussing on the starred chapters. `Recommended' here really does
mean \emph{recommended}; if you're struggling with any of the quizzes,
then looking through Sen's presentation of the formal work may very well
help.

In general (and I think this is true for every class), if you are doing
an essay on a topic, the recommended readings should be treated as
\emph{required}. If you choose to write on something, I expect you to be
on top of both the prose version of the topic, and the mathematical
version. In part two, this will be easy to follow, since Sen is very
good about separating out these versions.

The plan for the second half is not set in stone. This is a seminar
working through a book, and I think the best practice for classes like
that is to let the class itself determine the pace. I've made a best
guess at what speed we'll work at, but it's just that, a best guess. If
we need more time for something than I've guessed, we'll take more time.

The order will not change. There is a plan for how later parts build on
earlier parts, even within the Choose Your Own Adventure structure of
CCSW. If we don't get to stuff, it comes out of the end.

Changes to a week's reading will be posted on Canvas by the Friday
before. If nothing is posted, what is below stands.

Quizzes cover the material through the previous class meeting, whatever
that turned out to be, rather than whatever this schedule says it should
have been.

\subsection{Week 1: Decision and Strategic
Settings}\label{week-1-decision-and-strategic-settings}

\subsubsection{Tuesday, September 01 (Lecture
1)}\label{tuesday-september-01-lecture-1}

\textbf{Topic}: Course overview and the two organizing questions.
\textbf{Reading}: None.

\subsubsection{Thursday, September 03 (Lecture
2)}\label{thursday-september-03-lecture-2}

\textbf{Topic}: Decision theory under risk and uncertainty: expected
utility, Allais and Ellsberg. The formal contrast with strategic
settings.\\
\textbf{Reading}: Notes, 2.4-2.5 (Expected Value; Orthodox Decision
Theory).\\
\textbf{Recommended}: Notes, 2.1-2.3, if you want the probability
revision. Bonanno, 5.1-5.3, for a more careful statement of the von
Neumann-Morgenstern axioms.

\subsection{Week 2: Strategic Form
Games}\label{week-2-strategic-form-games}

\subsubsection{Tuesday, September 08 (Lecture
3)}\label{tuesday-september-08-lecture-3}

\textbf{Topic}: Strategic form games; dominance and iterated elimination
of dominated strategies.\\
\textbf{Reading}: Notes, Ch 1 (Basics of Game Theory).1.10.\\
\textbf{Recommended}: Bonanno, Ch 2 (Ordinal Games in Strategic Form),
2.1-2.2 and 2.5, which keeps the two iterated deletion procedures
carefully apart.

\subsubsection{Thursday, September 10 (Lecture
4)}\label{thursday-september-10-lecture-4}

\textbf{Topic}: Nash equilibrium.\\
\textbf{Reading}: Notes, 3.1-3.6.2 (Mixed Strategies through Finding
Best Responses).\\
\textbf{Recommended}: Bonanno, 2.6 (Nash equilibrium).

\textbf{Quiz}: Quiz 1 on Canvas, due Friday, September 11.

\subsection{Week 3: Mixed Strategies and
Rationalizability}\label{week-3-mixed-strategies-and-rationalizability}

\subsubsection{Tuesday, September 15 (Lecture
5)}\label{tuesday-september-15-lecture-5}

\textbf{Topic}: Mixed strategies and the interpretation problem (Aumann;
Harsanyi purification).\\
\textbf{Reading}: Notes, 3.6.3-3.8 (Finding Mixed Strategy Equilibria
through What Is a Mixed Strategy?).\\
\textbf{Recommended}: Bonanno, Ch 6 (Strategic-form Games), 6.1-6.3.

\subsubsection{Thursday, September 17 (Lecture
6)}\label{thursday-september-17-lecture-6}

\textbf{Topic}: Rationalizability; common knowledge of rationality.\\
\textbf{Reading}: Notes, 4.1 (Rationalizability).\\
\textbf{Recommended}: Bonanno, 6.4 (Strict dominance and
rationalizability), and Ch 10 (Rationality), 10.1-10.3. Ch 8 (Common
Knowledge) gives a more careful epistemic justification of
rationalizability.

\textbf{Quiz}: Quiz 2 on Canvas, due Friday, September 18.

\subsection{Week 4: Dynamic Games}\label{week-4-dynamic-games}

\subsubsection{Tuesday, September 22 (Lecture
7)}\label{tuesday-september-22-lecture-7}

\textbf{Topic}: Extensive form games; subgame perfection.\\
\textbf{Reading}: Notes, 5.1-5.4 (Normal Form and Extensive Form through
Subgame Perfect Equilibrium).\\
\textbf{Recommended}: Bonanno, 3.1-3.3; 4.1-4.4; and 7.2, which redoes
subgame perfection once payoffs are cardinal.

\subsubsection{Thursday, September 24 (Lecture
8)}\label{thursday-september-24-lecture-8}

\textbf{Topic}: Backward induction; the centipede paradox.\\
\textbf{Reading}: Notes, 5.5 (Problems with Backwards Induction).\\
\textbf{Recommended}: Bonanno, 3.2 and 3.4; 7.3 (Problems with the
notion of subgame-perfect equilibrium); 10.4 (Common knowledge of
rationality in extensive-form games).

\textbf{Quiz}: Quiz 3 on Canvas, due Friday, September 25.

\subsection{Week 5: Forward Induction and Bayesian
Games}\label{week-5-forward-induction-and-bayesian-games}

\subsubsection{Tuesday, September 29 (Lecture
9)}\label{tuesday-september-29-lecture-9}

\textbf{Topic}: Forward induction. Van Damme's burning-money game;
reasoning about what kind of player would have made an off-path move.\\
\textbf{Reading}: Notes, 5.6 (Money Burning Game).\\
\textbf{Recommended}: Elchanan Ben-Porath and Eddie Dekel, ``Signaling
future actions and the potential for sacrifice'', \emph{Journal of
Economic Theory} 57 (1992): 36-51.
\url{https://doi.org/10.1016/S0022-0531(05)80039-0}

\subsubsection{Thursday, October 01 (Lecture
10)}\label{thursday-october-01-lecture-10}

\textbf{Topic}: Bayesian games and Bayesian Nash equilibrium. Types,
incomplete information, and Nash equilibrium applied to type-conditional
strategies.\\
\textbf{Reading}: Notes, 6.1-6.5 (Two Kinds of Ignorance through
Purification).\\
\textbf{Recommended}: Bonanno, Ch 14 (Static Games), 14.1-14.3, and Ch
16 (The Type-Space Approach), 16.1-16.2. Also 9.5 (Harsanyi consistency
of beliefs), which is our common prior assumption under another name.

\subsection{Week 6: Signaling Games}\label{week-6-signaling-games}

\subsubsection{Tuesday, October 06 (Lecture
11)}\label{tuesday-october-06-lecture-11}

\textbf{Topic}: Signaling games; perfect Bayesian equilibrium.\\
\textbf{Reading}: Notes, 6.6-6.8 (Perfect Bayesian Equilibrium through
Signaling without Cooperation).\\
\textbf{Recommended}: Bonanno, 13.1-13.3, and 15.1 (One-sided incomplete
information).

\subsubsection{Thursday, October 08 (Lecture
12)}\label{thursday-october-08-lecture-12}

\textbf{Topic}: The beer-quiche game; the intuitive criterion.\\
\textbf{Reading}: Notes, 6.9 (The Intuitive Criterion). In-Koo Cho and
David M. Kreps, ``Signaling Games and Stable Equilibria'',
\emph{Quarterly Journal of Economics} 102 (1987): 179-221.
\url{https://www.jstor.org/stable/1885060}\\
\textbf{Recommended}: Bonanno, Ch 12 (Sequential Equilibrium). He does
not cover the intuitive criterion, but this is a similar idea.

\textbf{Quiz}: Quiz 4 on Canvas, due Friday, October 09.

\subsection{Week 7: Asymmetric
Information}\label{week-7-asymmetric-information}

\subsubsection{Tuesday, October 13 (Lecture
13)}\label{tuesday-october-13-lecture-13}

\textbf{Topic}: Spence on job-market signaling.\\
\textbf{Reading}: Michael Spence, ``Job Market Signaling'',
\emph{Quarterly Journal of Economics} 87 (1973): 355-374.
\url{https://www.jstor.org/stable/1882010}

\subsubsection{Thursday, October 15 (Lecture
14)}\label{thursday-october-15-lecture-14}

\textbf{Topic}: Akerlof on lemons.\\
\textbf{Reading}: George A. Akerlof, ``The Market for `Lemons': Quality
Uncertainty and the Market Mechanism'', \emph{Quarterly Journal of
Economics} 84 (1970): 488-500.
\url{https://www.jstor.org/stable/1879431}

\subsection{Week 8: Transition}\label{week-8-transition}

\subsubsection{Tuesday, October 20}\label{tuesday-october-20}

No class -- Fall Study Break (October 19-20).

\subsubsection{Thursday, October 22 (Lecture
15)}\label{thursday-october-22-lecture-15}

\textbf{Topic}: Introduction to CCSW.\\
\textbf{Reading}: CCSW: Ch 1 (Introduction); Ch A1 (Enlightenment and
Impossibility).\\
\textbf{Recommended}: CCSW: New Introduction (2017). This is a really
helpful overview about what the book is trying to do, and why it
includes \emph{these} topics.

\subsection{Week 9: Foundations}\label{week-9-foundations}

\subsubsection{Tuesday, October 27 (Lecture
16)}\label{tuesday-october-27-lecture-16}

\textbf{Topic}: The formal apparatus we'll use.\\
\textbf{Reading}: CCSW: Ch 2 (Unanimity); Ch 2* (Collective Choice Rules
and Pareto Comparisons), section 2\emph{1 only.\\
\textbf{Recommended}: CCSW: Ch 1} (Preference Relations). This is a
\emph{really} helpful overview of the terminology; you might want to
build a `cheat sheet' of key terms from it.

\subsubsection{Thursday, October 29 (Lecture
17)}\label{thursday-october-29-lecture-17}

\textbf{Topic}: Arrow's theorem, with a focus on why three options are
needed.\\
\textbf{Reading}: CCSW: Ch 3 (Collective Rationality).\\
\textbf{Recommended}: CCSW: Ch 3* (Social Welfare Functions); Ch 5
(Values and Choice); Ch 5* (Anonymity, Neutrality and Responsiveness).

\textbf{Essay}: First essay due Friday, October 30 at 5pm.

\subsection{Week 10: Arrow's Theorem, and the First
Escape}\label{week-10-arrows-theorem-and-the-first-escape}

\subsubsection{Tuesday, November 03 (Lecture
18)}\label{tuesday-november-03-lecture-18}

\textbf{Topic}: Sen's proof of Arrow's Theorem.\\
\textbf{Reading}: CCSW: Ch A1* (Social Preference), to p.~288.\\
\textbf{Recommended}: CCSW: Ch 3* (Social Welfare Functions).

\subsubsection{Thursday, November 05 (Lecture
19)}\label{thursday-november-05-lecture-19}

\textbf{Topic}: Views which reject Universal Domain, in particular the
single-peakedness assumption\\
\textbf{Reading}: CCSW: Ch 10 (Majority Choice and Related Systems).\\
\textbf{Recommended}: CCSW: Ch 10* (Restricted Preferences and Rational
Choice).

\textbf{Quiz}: Quiz 5, on Canvas, due Friday, November 06.

\subsection{Week 11: Two More Escapes}\label{week-11-two-more-escapes}

\subsubsection{Tuesday, November 10 (Lecture
20)}\label{tuesday-november-10-lecture-20}

\textbf{Topic}: Views which reject Collective Rationality, with a focus
on Gibbard's oligarchy result.\\
\textbf{Reading}: CCSW: Ch A1* (Social Preference), pp.~289-293; Ch 4
(Choice Versus Orderings); Ch A2 (Rationality and Consistency).\\
\textbf{Recommended}: CCSW: Ch 4* (Social Decision Functions); Ch A2*
(Problems of Social Choice).

\subsubsection{Thursday, November 12 (Lecture
21)}\label{thursday-november-12-lecture-21}

\textbf{Topic}: Third escape: use more than preferences.\\
\textbf{Reading}: CCSW: Ch 7 (Interpersonal Aggregation and
Comparability).\\
\textbf{Recommended}: CCSW: Ch 7* (Aggregation Quasi-Orderings); Ch 8
(Cardinality With or Without Comparability).

\textbf{Quiz}: Quiz 6, on Canvas, due Friday, November 13.

\subsection{Week 12: Strategy and
Liberty}\label{week-12-strategy-and-liberty}

\subsubsection{Tuesday, November 17 (Lecture
22)}\label{tuesday-november-17-lecture-22}

\textbf{Topic}: Strategic manipulation, up to the Gibbard-Satterthwaite
theorem. (This brings game theory back into social choice theory.)\\
\textbf{Reading}: Amartya Sen, ``The Possibility of Social Choice'',
\emph{American Economic Review} 89 (1999): 349-378.
\url{https://www.aeaweb.org/articles?id=10.1257/aer.89.3.349}\\
\textbf{Recommended}: Allan Gibbard, ``Manipulation of Voting Schemes: A
General Result'', \emph{Econometrica} 41 (1973): 587-601.
\url{https://www.jstor.org/stable/1914083}

\subsubsection{Thursday, November 19 (Lecture
23)}\label{thursday-november-19-lecture-23}

\textbf{Topic}: Sen's liberal paradox.\\
\textbf{Reading}: CCSW: Ch 6 (Conflicts and Dilemmas); Ch 6* (The
Liberal Paradox).\\
\textbf{Recommended}: Amartya Sen, ``The Impossibility of a Paretian
Liberal'', \emph{Journal of Political Economy} 78 (1970): 152-157.
\url{https://doi.org/10.1086/259614}

\textbf{Quiz}: Quiz 7, on Canvas, due Friday, November 20.

\subsection{Week 13: Rights}\label{week-13-rights}

\subsubsection{Tuesday, November 24 (Lecture
24)}\label{tuesday-november-24-lecture-24}

\textbf{Topic}: Rights and social choice. Sen's later response to his
own paradox.\\
\textbf{Reading}: CCSW: Ch A5 (The Idea of Rights).\\
\textbf{Recommended}: CCSW: Ch A5* (Rights and Social Choice).

\subsubsection{Thursday, November 26}\label{thursday-november-26}

No class -- Thanksgiving recess (November 25-27).

\subsection{Week 14: Equity, Justice,
Capability}\label{week-14-equity-justice-capability}

\subsubsection{Tuesday, December 01 (Lecture
25)}\label{tuesday-december-01-lecture-25}

\textbf{Topic}: Equity and justice. The weak equity axiom, and the
limits of welfarism.\\
\textbf{Reading}: CCSW: Ch 9 (Equity and Justice).\\
\textbf{Recommended}: CCSW: Ch 9* (Impersonality and Collective
Quasi-Orderings); Ch A3 (Justice and Equity).

\subsubsection{Thursday, December 03 (Lecture
26)}\label{thursday-december-03-lecture-26}

\textbf{Topic}: The capability approach. \textbf{Reading}: Amartya Sen,
``Equality of What?'', Tanner Lecture on Human Values, Stanford
University, 22 May 1979; reprinted in \emph{Choice, Welfare and
Measurement}.
\url{https://tannerlectures.org/lectures/equality-of-what/}\\
\textbf{Recommended}: CCSW: Ch A3* (Social Welfare Evaluation). Martha
Nussbaum, ``Capabilities as Fundamental Entitlements: Sen and Social
Justice'', \emph{Feminist Economics} 9 (2003): 33-59.
\url{https://doi.org/10.1080/1354570022000077926}

\textbf{Quiz}: Quiz 8, on Canvas, due Friday, December 04.

\subsection{Week 15: Democracy and Public
Reasoning}\label{week-15-democracy-and-public-reasoning}

\subsubsection{Tuesday, December 08 (Lecture
27)}\label{tuesday-december-08-lecture-27}

\textbf{Topic}: Democracy and public engagement. \textbf{Reading}: CCSW:
Ch A4 (Democracy and Public Engagement).\\
\textbf{Recommended}: CCSW: Ch A4* (Votes and Majorities).

\subsubsection{Thursday, December 10 (Lecture
28)}\label{thursday-december-10-lecture-28}

\textbf{Topic}: Reasoning and social decisions. \textbf{Reading}: CCSW:
Ch A6 (Reasoning and Social Decisions).\\
\textbf{Recommended}: CCSW: Ch 11 (Theory and Practice).

\textbf{FINAL ESSAY}: Due Friday, December 18 at 5pm.

\newpage

\section{Plagiarism}\label{plagiarism}

Although team-work, and even co-authorship, is encouraged, plagiarism is
strictly prohibited. You are responsible for making sure that none of
your work is plagiarized. Be sure to cite work that you use, both direct
quotations and paraphrased ideas. Any citation method that is tolerably
clear is permitted, but if you'd like a good description of a citation
scheme that works well in philosophy, look at
\url{http://bit.ly/VDhRJ4}.

You are encouraged to discuss the course material, including
assignments, with your classmates, but all written work that you hand in
under your own name must be your own. If work is handed is as the work
of two people, you are affirming that each person did a fair share of
the work. (Note that when you're submitting work on Canvas, you have to
each submit the paper, even if it is co-authored. That way Canvas knows
that everyone has turned in work.)

This is a 400-level course, so I don't want to be as heavy-handed with
enforcement as in a big lecture course. But still we have some basic
ground rules.

\begin{enumerate}
\def\labelenumi{\arabic{enumi}.}
\tightlist
\item
  Acknowledge all assistance that you got. Note that this is something
  Sen does somewhat briefly in the original Preface to CCSW, and at much
  greater length in the revised version.
\item
  Don't get any assistance from an LLM that you wouldn't be happy to get
  from a human. List precisely what help you do get from LLMs, and keep
  a record of all interactions. (Don't use anyone else's account.)
\item
  When I say the work is `your own', I mean that if someone asked you
  about any sentence in it, you could answer two questions: \emph{What
  does this mean?} and \emph{Why did you say this?}. If I suspect you
  are turning in work that is not your own, I'll be very disappointed,
  and the first step will be to ask you those two questions.
\end{enumerate}

You should also be familiar with the academic integrity policies of the
College of Literature, Science \& the Arts at the University of
Michigan, which are available here:
\url{http://www.lsa.umich.edu/academicintegrity/}. Violations of these
policies will be reported to the Office of the Assistant Dean for
Student Academic Affairs, and sanctioned with a course grade of F.

\section{Disability}\label{disability}

The University of Michigan abides by the Americans with Disabilities Act
of 1990, Section 504 of the Rehabilitation Act of 1973, and other
applicable federal and state laws that prohibit discrimination on the
basis of disability, which mandate that reasonable accommodations be
provided for qualified students with disabilities.

If you have a disability, and may require some type of instructional
and/or examination accommodation, please contact me early in the
semester. If you have not already done so, you will also need to
register with the Office of Services for Students with Disabilities. The
office is located at G664 Haven Hall.

For more information on disability services at the University of
Michigan, go to \url{http://ssd.umich.edu}.




\end{document}
